\section{Related work} \label{sec:related}

Two families of methods dominate the prediction of solid solubility and solvation, and they make
opposite choices about physical structure. Physics-informed pipelines such as SolProp
\cite{vermeire2022solprop} assemble a thermodynamic cycle from separately learned components
(solvation free energy, vapour pressure, fusion properties), so that each prediction is decomposable
and its intermediate quantities carry physical meaning. Direct models forgo that structure: FastSolv
\cite{osanloo2024} regresses solubility from molecular descriptors with a deep neural network and is the
current leader on extrapolation to new solutes, a task on which inter-laboratory disagreement sets a
substantial irreducible-error floor. On accuracy the ordering is clear: the black box leads and the
physics-informed cycle trails it.

The activity coefficient is the quantity a solubility model must supply, and learned surrogates for it
have advanced quickly. Winter et al.\ \cite{winter2022} predict infinite-dilution activity coefficients
\lngi\ directly from SMILES with a language model, and Sanchez-Medina et al.\ \cite{sanchezmedina2023}
embed a Gibbs--Helmholtz constraint in a graph network to capture their temperature dependence. These
approaches treat the activity coefficient as the regression target, rather than retaining a mechanistic
model that consumes a predicted molecular property.

The broader question of when a fixed physical structure helps or hurts a learned model is central to
physics-informed machine learning \cite{karniadakis2021}. Its dominant forms impose structure as a
differentiable object: physics-informed neural networks enforce a governing equation through the loss
\cite{raissi2019}, while physics-guided and gray-box hybrids compose learned components with fixed
mechanistic maps \cite{willard2022}, and a recent operator-learning approach adds a trainable correction
to a misspecified governing equation \cite{ma2026}. The prevailing expectation is that more physics
improves generalisation, yet the same structure can degrade it: physics-informed networks have
documented failure modes in which the imposed constraint makes optimisation or accuracy worse
\cite{krishnapriyan2021}. In neural solvers of partial differential equations, Mohan et al.\
\cite{mohan2024} show that a learnable term added to a fixed discretised operator can absorb the
simulation's truncation artifacts rather than the physics, so that the program's apparent
interpretability is illusory; their error source is numerical, and they have no external reference for
the learned term.

The same concern appears in concept-bottleneck models as \emph{concept leakage}, in which a learned
concept silently encodes information beyond its nominal meaning and its interpretation becomes misleading
\cite{mahinpei2021,margeloiu2021,havasi2022}. Recent work traces one route to leakage to the downstream
model rather than to a trainable head: an inflexible or misspecified head forces the encoder to deform
the concept into carrying a dependence the head cannot represent \cite{parisini2025}, and overwriting a
frozen head's leaked concepts with their ground-truth values can then degrade accuracy
\cite{espinosa2025leakage}.

Several diagnostic tools are relevant to such failures. The gap between a fixed physical map and reality
is the model-discrepancy problem of Kennedy and O'Hagan \cite{kennedy2001} and its calibration critique
\cite{brynjarsdottir2014}; separating whether predictions are right on average from whether they are
sharp is the calibration and reliability--resolution decomposition of Murphy and DeGroot
\cite{murphy1973,degroot1983}; and a misspecified model still converges to a well-defined pseudo-true
limit under the quasi-maximum-likelihood account of White \cite{white1982}. That a decomposition may be
only weakly constrained by the data is the subject of parameter sloppiness
\cite{transtrum2015,gutenkunst2007} and of structural versus practical identifiability
\cite{raue2009,villaverde2016}, and the risk that apparent skill reflects pipeline choices rather than
physics is the concern of underspecification \cite{damour2022}, leakage and reproducibility
\cite{kapoor2023}, and compositional risk \cite{mahajan2025crm}.

What this literature lacks is a way to \emph{measure}, rather than merely diagnose, whether a fixed
physical model or its learned inputs limits a composed predictor's accuracy, and to test whether the
learned intermediate has become a surrogate that compensates for the model's error. The concept-leakage
literature identifies the phenomenon but has no external reference against which to quantify it, and the
neural-solver work has a fixed operator but only a numerical, unreferenced error term. We supply the
missing ingredient: an external reference for the learned intermediate---a reference $\sigma$-profile
for solubility, and tabulated substituent constants for a second chemical domain---which turns the
qualitative concern into a model-free measurement that separates model misspecification from input
insufficiency. Using it, we show on a fixed thermodynamic model that the learned intermediate becomes a
low-rank, transferable surrogate, and that grounding it on reference values can therefore reduce
accuracy. Our data are drawn from BigSolDB~2.0 \cite{krasnov2025bigsoldb}; FastSolv and SolProp remain
the accuracy reference points, which we do not contest.
